\documentclass[../../../include/open-logic-section]{subfiles}

\begin{document}

\olfileid{sth}{story}{urelements}
\olsection{是否包含本元？}

在接下来的几章中，我们将尝试从集合的累积迭代概念中提取公理。但在进一步讨论之前，我们需要对\emph{本元}多作一些说明。

\olref[approach]{sec}节的图景只允许我们由已有的\emph{集合}来形成新的集合。然而，我们可能希望允许某些\emph{非集合}——牛、猪、沙粒或其他任何事物——成为集合的元素。如果是这样，我们将从某些基本元素（即\emph{本元}）出发，然后说在每一阶段~$S$，我们会形成由本元或在~$S$ 之前的阶段形成的集合（以任意方式组合）所组成的“所有可能的”集合。最终得到的图景将更像这样：

\begin{center}
	\begin{tikzpicture}[scale=0.6]
	\tikzset{cut_here/.style={densely dotted}}
	\draw (-4,6) -- (-1, 0) -- (1,0) -- (4,6); 
	\draw[cut_here] (-5,8)--(-4,6);
	\draw[cut_here] (5,8)--(4,6);
	\draw(-1.5, 1)--(1.5, 1);
	\draw(-2, 2)--(2, 2);
	\draw(-2.5, 3)--(2.5,3);
	\draw(-3, 4)--(3, 4);
	\draw(-3.5, 5)--(3.5, 5);
	\draw(-4, 6)--(4, 6);
	\node[label] at (2, 0) {\small 0};
	\node[label] at (2.5, 1) {\small 1};
	\node[label] at (3, 2) {\small 2};
	\node[label] at (3.5, 3) {\small 3};
	\node[label] at (4, 4) {\small 4};
	\node[label] at (4.5, 5) {\small 5};
	\node[label] at (5, 6) {\small 6};
	\end{tikzpicture}
\end{center}

所以现在我们需要做一个决定：\emph{我们应该允许本元吗？}

从哲学上讲，将本元纳入我们的理论建构是有道理的。这样做的主要理由是使我们的集合论具有\emph{适用性}。为了说明这一点，回顾\olref[sfr][siz][]{chap}，我们说两个集合~$A$ 和~$B$ 大小相同，即 $\cardeq{A}{B}$，当且仅当它们之间存在一个双射。现在，如果田野里的牛和猪圈里的猪都能构成集合，我们就可以对“牛和猪一样多”这一论断提供一种集合论的处理方式。但如果我们禁止本元，使得牛和猪\emph{不}构成集合，那么这种集合论的处理方式就不可用了。事实上，我们将无法直接将集合论应用于集合本身以外的任何事物。（关于包含本元的更多理由，参见 \citealt[pp.~vi, 24, 50--1]{Potter2004}。）

然而，在数学上，允许本元的情况相当罕见。部分原因是，表述不含本元的集合论要\emph{略微}容易一些。但偶尔，人们也会给出更有趣的理由来将本元排除在集合论之外：

\begin{quote}
	根据集合论作为数学基础这一信念，我们应该能够仅通过谈论集合来涵盖所有数学，因此我们的变元不应遍历像牛和猪这样的对象。
	%But if $C$ is a cow, $\{C\}$ is a set, but not a legitimate mathematical object. 
	\citep[p.~8]{Kunen1980}
\end{quote}

因此：关注适用性倾向于\emph{包含}本元；而关注还原论的基础目标（将数学还原为纯集合论）则可能倾向于\emph{排除}它们。一点轻微的惰性，也指向了排除本元的方向。

我们将选择最省力的路径。不过，这在一定程度上也有教学法上的理由。我们的目标是向你介绍入门集合论哲学所需的集合论基础知识。而大部分相关的哲学文献讨论的都是\emph{不}包含本元的形式化集合论。因此，如果本书能相当贴近那些文献，或许会更有用。

\end{document}